% page 454 du fac-simile (image p-453 du scan)
% bloc 157.5 x 237.7 mm ; marge gauche 8.0 mm
\pgc{-12.700mm}{0.000mm}{
\l{\cel{3}446}
\l{}
\l{\sou{pledar}. (\sou{netrans.}) Defensar (afero procesala) koram la ju-}
\l{\cel{3}diciisti. - DEFS.}
\l{}
\l{\sou{plektar}. (\sou{trans.}) Fasonar texuro (ek fasko de hari, e c.)}
\l{\cel{3}plata, rubando-forma, per interkrucumo. - EL.}
\l{}
\l{\sou{plektro.}\cel{2}(\sou{epoki antiqua}). Mikra vergo ek ivoro, quan onu}
\l{\cel{3}uzis por igar vibrar la kordi di liro. - DEFIS.}
\l{}
\l{\sou{plena}. (de) I. Qua kontenas la tota quanto quan lu povas}
\l{\cel{3}kontenar. - II. Qua kontenas granda quanto (kom ex. :}
\l{\cel{3}plenigar ye ta o ca kozo,\cel{1}\sou{per}\cel{1}ta o ca instrumento). III.}
\l{\cel{3}(\sou{fiz.}) Plenigita tote, ye materio solida, liquida, o gasa.}
\l{\cel{3}EFIS.}
\l{}
\l{\sou{plendar}.\cel{2}(\sou{netrans., pri)}. I. Expresar sua deskontenteso pri}
\l{\cel{3}ulu od ulo. - II. Expresar, che autoritatozo policala, sua}
\l{\cel{3}motivo di deskontenteso por obtenar protekto, indemno, e c.}
\l{\cel{3}- EFS.}
\l{}
\l{\sou{plenipotenco}. Yuro quan onu recevas de ulu pri agar ye lua}
\l{\cel{3}nomo. (Dicesas precipue pri la yuro donita da suvereno, pri}
\l{\cel{3}agar ye guvernistaro o korto stranjera.) - DEFIS.}
\l{}
\l{\sou{pleomorfa}. (\sou{mikrobiol.})\cel{2}Di qua la formo varias efike daula}
\l{\cel{3}kondicioni determinita. -}
\l{}
\l{\sou{pleonasmo.}\cel{2}I. (\sou{gram.)}\cel{2}Uzo di vorti qui esas superflua pri}
\l{\cel{3}la senco. - II. (\sou{retor.}) Abundego di la termini por insis-}
\l{\cel{3}tar pri la penso. (kom ex.: \textquotedbl{}Me vidis lu, ya, vidis lu per}
\l{\cel{3}mea okuli propra\textquotedbl{}. - DEFIRS.}
\l{}
\l{\sou{plesiosauro}. (\sou{paleont.})\cel{2}Genero di reptero, tipo di la familio}
\l{\cel{3}\textquotedbl{}plesiosauridei\textquotedbl{}, grandanimalo marala, 3/5 metri longa, re-}
\l{\cel{3}markinda pro lua mikra lezardo-kapo, lua denti strioza, lua}
\l{\cel{3}kolo longa nemezureble. - L.\cel{1}\sou{plesiosaurus}. - DEFIS.}
\l{}
\l{\sou{pleto}. Pladego de metalo, lako, porcelano, kristalo,sur qua}
\l{\cel{3}onu servas teo, kafeo, o cetera drinkaji, e sur qua onu}
\l{\cel{3}pozas la bako-kukaji. - F.}
\l{}
\l{\sou{pletoro.}\cel{1}I. (\sou{patol.)}\cel{2}Pleneso di la vaskuli pro abundego di}
\l{\cel{3}sango, di humori. - II. (\sou{metaf.}) Tromulteso quaze morba.}
\l{\cel{3}- DEFIS.}
\l{}
\l{\sou{pleuronekt}o. (\sou{zool}.) Marfisho plata, di qua la korpo esas rom-}
\l{\cel{3}boida, analoga a turboto, ma plu ovala. - L.\cel{1}\sou{pleuronectes}.L}
\l{}
\l{\sou{pleuronektoidi}. (\sou{zool.)}\cel{2}Genero de fishi plata qui natas sur}
\l{\cel{3}sua flanko. L.\cel{1}\sou{pleuronectidae}. - L.}
\l{}
\l{plevro. (anat.) Membrano seroza, qua tegas la pulmoni e kovras}
\l{\cel{3}la parieti di la torako. - EFIRS.}
}
